% Argumenten las principales diferencias observadas en las dinámicas del pozo nito simétrico y no simétrico. Les recomiendo que varíen los parámetros y analicen diferentes situaciones que puedan ocurrir.

% ¿Cuáles conceptos aprendimos con este proyecto? ¿Cuáles habilidades desarrollamos con el desarrollo del proyecto?

\subsection{Andrés Saul Pérez Martínez }
a primera diferencia en los dos casos son las soluciones que nos proporciona la simetría, con respecto a la paridad de las soluciones que podamos obtener. Esto afecta a la dinámica principalmente en las fronteras, donde al ser simétricas los rebotes son parecidos y respeta cierta simetría; en cambio, al tener potenciales diferentes en el no simétrico podemos observar que estos cambios generan que en las fronteras la onda varíe y cambie el comportamiento. Por ejemplo, el pulso gaussiano se ve afectado por esta diferencia de paredes al momento de que va variando en el tiempo gracias a esto, lo que en consecuencia genera que haya momentos en donde el momento prefiera estar donde hay un menor potencial de pared que del otro lado, ya que si es más probable encontrarlo por ahí mismo por la interferencia generada, lo cual es diferente para el simétrico.

Con este proyecto retomé la programación y, a pesar que en el semestre ya había podido reforzarme, la verdad es que me di cuenta que cuando ya es algo más largo se te van acabando las variables y luego sale error porque ya estaba todo mezclado. Entonces pude desarrollar el poder aprender mejor a programar en Julia y leer documentación o reconocer en qué partes se podía usar funciones que no conocía, propias o librerías, para por ejemplo raíces más rápidas o con funciones analíticas. Del caso de cuántica me gustó poder aterrizar las ideas cuando ya tenemos en cuenta la variable del tiempo, porque al inicio del pozo simétrico yo estaba convencido si era continuo y con esto me doy cuenta si era una especie de discreto y aterrizo varios conceptos.